The electrode-level tests above establish which individual electrodes were affected by operating their region independently. The region-mean analysis that follows instead aggregates those same electrode results to ask whether the region as a whole changed relative to the same region within the full montage. Electrode-level results were aggregated into a single region-mean $F_1$ per subset, following the aggregation used for Table~\ref{tab:region_performance} and Figure~\ref{fig:region_performance}. Each electrode's own macro $F_1$ was first averaged over sessions, and those per-electrode means were then averaged across the electrodes the subset contained, applied here to the subset's own self-contained rerun rather than to the full-montage run. Each subset's region-mean $F_1$ was compared against the same region's electrodes scored inside the full-montage run, on the same sessions, rather than against the full-montage headline oracle $F_1$, using the same paired-testing framework as the electrode-level comparison above, a two-tailed Wilcoxon signed-rank test paired by session with the matched-pairs rank-biserial correlation as effect size, now Bonferroni-corrected over the 13 subsets compared within each dataset rather than over individual electrodes. Because Stages~A--C were rerun independently for each subset, this design measured whether isolating a region as its own detector changed how well that region's own electrodes performed. It was not a leave-one-region-out analysis and did not establish that any region was necessary to the complete system.

The region-mean comparisons showed that every anatomical subset, in both datasets, scored significantly lower than the same electrodes' full-montage region-mean $F_1$ after Bonferroni correction, with no subset in either corpus failing to reach significance (Table~\ref{tab:exp1_subset_summary}, Figure~\ref{fig:exp1_region_boxplot}). The frontal subset was the strongest-performing region-mean comparison in both corpora. On Internal, the frontal subset reached $F_1=66.92\%$ against $F_1=69.68\%$ for the same six electrodes scored inside the full montage, a paired difference of $-2.76$ percentage points ($p<0.001$, $r=-0.80$). On Cao2018, the frontal subset reached $F_1=66.61\%$ against $F_1=68.91\%$ for the same four electrodes, a paired difference of $-2.29$ percentage points ($p=0.026$, $r=-0.47$). Both differences were statistically significant, but modest in absolute terms relative to the region's own $F_1$. The Cao2018 frontal comparison additionally had the largest corrected $p$-value and the smallest effect size of any coarse region reported below.

The remaining regions showed the same shortfall pattern at much lower absolute performance levels. On Internal, the central, parietal, occipital and posterior subsets reached region-mean $F_1$ of $13.99\%$, $1.72\%$, $2.16\%$ and $2.08\%$, respectively, against $14.80\%$, $1.94\%$, $2.30\%$ and $2.18\%$ for the same electrodes in the full montage, paired differences of $-0.81$, $-0.22$, $-0.15$ and $-0.11$ percentage points (all $p<0.001$). On Cao2018, the corresponding subsets reached $34.29\%$, $15.00\%$, $6.42\%$ and $8.62\%$, against $38.13\%$, $18.70\%$, $8.24\%$ and $10.86\%$ for the full-montage electrodes, paired differences of $-3.83$, $-3.69$, $-1.83$ and $-2.24$ percentage points (all $p<0.001$).

Two conclusions follow from these region-mean comparisons, both drawn from electrode-level and region-mean comparisons in which every electrode retained its dataset's full session pool. First, the spatial hierarchy established in the full-montage analysis was preserved within the self-contained subset reruns. Frontal electrodes clearly outperformed the central, parietal, occipital and posterior regions in both datasets, exactly as they did within the full 32-channel run. Second, operating an anatomical region as its own self-contained detector generally reduced the performance of that region's own electrodes relative to their performance inside the full montage. Every paired difference was negative in both corpora, the Internal frontal decrease ($-2.76$ points) was larger than any other Internal region's decrease, and the Cao2018 frontal decrease ($-2.29$ points) was smaller than the central ($-3.83$ points) and parietal ($-3.69$ points) decreases and close to the posterior decrease ($-2.24$ points). Matched-pairs rank-biserial correlations were substantial for every region-mean comparison ($|r|\geq0.65$) except the Cao2018 frontal subset, where the effect was smaller ($r=-0.47$). The direction of the effect, subset $F_1$ below montage $F_1$, was nonetheless consistent across every comparison in both datasets.
